The discrete-state regime-switching tradition opens with Hamilton~\cite{hamilton1989new}, who introduced Markov-switching to economics. Ryd\'{e}n et al.~\cite{ryden1998stylized} fitted Gaussian-emission HMMs at $K = 2$ or $3$ zero-mean states to subseries of the 1928--1991 S\&P~500 index and concluded that the slow decay of the absolute-return ACF was the one stylized fact the model could not generate by maximum likelihood. Bulla and Bulla~\cite{bulla2006stylized} recovered the fit with hidden semi-Markov models, replacing the geometric sojourn with an explicit-duration family at the cost of a duration-augmented forward--backward, and Nystrup et al.~\cite{nystrup2017long} extended the same line to long-memory volatility. The per-state densities used here draw on the maximum-likelihood Student-$t$ formulations of Peel and McLachlan~\cite{peel2000robust} and Liu and Rubin~\cite{liu1995ml}, with the generalized error distribution (GED) providing a continuous shape parameter that interpolates between Gaussian ($p = 2$) and Laplace ($p = 1$). The present paper contributes a common estimation and evaluation framework across these families rather than a new model class, and tests empirically whether the finite collection of eigenvalue-driven ACF modes is adequate on the studied data rather than claiming general long-memory decay.

Conditional-variance models form a parallel line in which the latent state is the volatility process rather than a discrete regime: Gaussian GARCH~\cite{engle1982autoregressive, bollerslev1986generalized}, heavy-tailed innovations~\cite{bollerslev1987conditionally}, and Markov-switching GARCH (MS-GARCH)~\cite{haas2004new}, with the \texttt{MSGARCH} package of Ardia et al.~\cite{ardia2019msgarch} the reference implementation. These classes can produce heavy-tailed predictive distributions and conditional VaR forecasts; unimodality alone imposes no excess-kurtosis ceiling. The distinction relevant here is narrower: the tested CHMM assigns a separate location, scale, and potentially shape parameter to each finite state, whereas the tested conditional-volatility baselines place most regime dependence in the variance recursion, and the empirical comparison separates those specifications on marginal fit and excess kurtosis without establishing structural impossibility for the broader families.

Deep generative market models offer high-capacity distributional and temporal representations, often with less direct state-level interpretation. GAN-based generators require temporal structure in the architecture or training objective to reproduce volatility clustering~\cite{takahashi2019modeling}; prominent examples include TimeGAN~\cite{yoon2019timegan} and the convolutional Wasserstein GAN QuantGAN~\cite{wiese2020quantgan}. We include one in-house QuantGAN implementation---a materially smaller approximation of the reference architecture, not a faithful reproduction---as the deep-generative row of the main panel, read as a negative control on whether that architecture recovers the stylized facts at our sample size rather than as a representative baseline for the whole family; reports of tail mode collapse in GAN-based financial generators~\cite{takahashi2019modeling} motivate the heavy-tail diagnostic but do not establish that failure for every deep formulation. Synthetic-data evaluation is anchored on the framework of Stenger et al.~\cite{stenger2024thinking} and proper scoring rules~\cite{gneiting2007strictly}; the walk-forward evaluation and periodic-refit recommendation draw on the structural-break and online-updating literature~\cite{cappe2011online}.
